\documentclass[11pt,a4paper]{article}
\usepackage{fontspec}
\usepackage{amsmath,amssymb,amsthm}
\usepackage[margin=1in]{geometry}
\usepackage{longtable,booktabs,array,calc}
\usepackage{graphicx}
\usepackage{float}
\floatplacement{figure}{H}
\usepackage{needspace}
\allowdisplaybreaks[1]
\clubpenalty=10000
\widowpenalty=10000
\raggedbottom
\usepackage{fvextra}
\DefineVerbatimEnvironment{verbatim}{Verbatim}{breaklines=true,fontsize=\small}
\usepackage{xurl}
\usepackage[hidelinks]{hyperref}
\hypersetup{pdftitle={The Juggler Map and the 3n±1 Maps: Exact Coding
and Arithmetic Obstructions},pdfauthor={Philippe
Cochin},pdfsubject={Mathematics preprint}}
\urlstyle{tt}
\setlength{\emergencystretch}{3em}
\setlength{\LTpre}{1em}
\setlength{\LTpost}{1em}
\setlength{\tabcolsep}{4pt}
\renewcommand{\arraystretch}{1.12}
\providecommand{\tightlist}{\setlength{\itemsep}{0pt}\setlength{\parskip}{0pt}}
\providecommand{\passthrough}[1]{#1}
\setcounter{secnumdepth}{-1}
\title{The Juggler Map and the 3n±1 Maps\\[0.5em]{\large Exact Coding
and Arithmetic Obstructions}}
\author{Philippe Cochin \\ \texttt{philippe@cochin.fr} \\
  \small ORCID \href{https://orcid.org/0009-0004-1939-3382}{0009-0004-1939-3382}}
\date{25 September 2026 · Version 0.8.1}
\begin{document}
\maketitle
\begin{abstract}

The Juggler map and the shortcut \(3n+1\) and \(3n-1\) maps share the
word multiplier \(3^o/2^L\), but realize their words in different
arithmetic and metric spaces. Classical parity coding defines an exact
map from every Juggler orbit into the 2-adic \(3n-1\) system. On actual
periodic Juggler orbits it preserves every return time, and its value is
an ordinary integer exactly when an explicit word divisibility holds. A
family of genuine Juggler prefixes shows that arbitrarily precise
modular return does not force that divisibility. We count this
constructed family with an exact leading constant and prescribe its
residues throughout the final even run. For the fixed word OOE, we also
give an explicit counting error and a uniform polynomial first-witness
bound in the modulus; this quantitative extension has a complete Lean
proof, with independent review pending. On the positive integers, we
prove that every \(3n-1\) target prime to three has at least
\(X^{21/25}\) ancestors below \(X\), for all sufficiently large \(X\).
This signed adaptation uses height-corrected inverse-tree inequalities,
a closed root domain, and an exact integer certificate with 177147 rows.
A mean inequality places every positive certificate in the fixed
\(1/50\) grid, for either sign and at every finite residue level,
strictly below exponent \(0.99\) in the at-most-linear range. For
reciprocal mass on the signed Collatz fibres, no finite ternary weight
table reproduces itself for either sign, over any fixed number of
generations or under any bounded stopping rule; divergent ancestor mass
at a nonperiodic target follows from divergence of an explicit
coefficient series, a premise that remains open. Lean formalizations
accompany the orbit, period, counting, grid-ceiling, modular-return, and
signed-fibre results, including the modular return's analytic recurrence
input. These results separate symbolic correspondence from integer
realization and reciprocal-mass transport; no universal termination or
unrestricted cycle exclusion theorem is asserted.

\end{abstract}

\textbf{2020 Mathematics Subject Classification:} Primary 11B83;
Secondary 11K31, 11Y16, 68V15.

\textbf{Keywords:} Juggler map; signed Collatz maps; 2-adic parity
coding; inverse trees; computer-assisted proof; Lean.

\section{1. Introduction}\label{introduction}

We compare the maps \[
J(n)=
\begin{cases}
\lfloor\sqrt n\rfloor,&n\ \text{even},\\
\lfloor n^{3/2}\rfloor,&n\ \text{odd},
\end{cases}
\qquad
C_\sigma(x)=
\begin{cases}
x/2,&x\ \text{even},\\
(3x+\sigma)/2,&x\ \text{odd},
\end{cases}
\quad \sigma\in\{-1,+1\}.
\] Unless a larger domain is specified, inputs are positive integers.
Write \(g=C_{-1}\). On ordinary integers \(C_{+1}(-x)=-g(x)\); positive
\(3n-1\) orbits correspond to negative \(3n+1\) orbits, not to positive
\(3n+1\) orbits.

Along a prescribed branch word, the leading Collatz multiplier acts on
\(x\), whereas the leading Juggler multiplier acts on \(\log n\). The
same letters generate the same formal walk one logarithm apart. Turning
this observation into a statement about integer trajectories requires
additional information. This paper asks which information survives the
correspondence.

Sections 2 and 3 construct the exact code and prove period preservation.
Section 4 exhibits actual modular returns whose associated periodic-word
codes are nonintegral. Section 5 proves a signed ancestor-count theorem,
including the height correction that prevents an automatic transfer of
the positive-map proof. Section 6 proves a limitation of the fixed grid
used in that argument. Section 7 distinguishes counting from reciprocal
mass and bounded word identities from unbounded stopping identities, and
reduces reciprocal mass on the signed Collatz fibres to one coefficient
series.

\subsection{1.1. Contributions and prior
work}\label{contributions-and-prior-work}

Terras {[}T76{]} supplies the classical parity-vector bijection and
stopping-word combinatorics. Bernstein and Lagarias {[}BL96{]} supply
the 2-adic inverse parity map, including the signed affine family. We
use that construction, rather than claim a new coding mechanism. Prasad
and Prasad {[}PP25{]} discuss the leading-word approximation for
Juggler-like sequences. The shared multiplier and the word counts are
background.

The Juggler-specific refinement in Section 3 is preservation of every
return time on an actual periodic orbit. Its proof combines classical
coding with monotonicity within each branch. The integrality criterion
is the classical affine cycle equation stated alongside this refinement.
Section 4 gives a quantified obstruction involving actual floor
iterates. Its counting and residue corollaries extract further
consequences of classical power equidistribution; no new general
distribution criterion is claimed. Theorem 4.4 adds a quantitative
specialization for the single word OOE.

Krasikov and Lagarias {[}KL03, Theorem 6.1{]} prove an ancestor exponent
\(0.84\) for every positive target prime to three under \(C_{+1}\).
Their theorem is not a statement about positive \(g\)-targets. Our
Section 5 proves the latter statement by a strict-grid construction. The
\(1/50\) cap table and induction measure follow M. Sharpe's formal
implementation {[}S26{]}, with attribution and its MIT notice retained
in the source. The signed height estimates, the domain for arbitrary
positive unit targets, and the independently generated certificate are
supplied here. We claim neither priority for the grid method nor a
larger positive-map exponent.

Section 6 isolates a common mean obstruction for this fixed grid at
every finite residue level. This is a limitation of those certificate
inequalities, not a bound on actual ancestor growth.

The complete odd-predecessor fibre of a Collatz target is one ray
\(n\mapsto4n+s\) {[}Z08{]}, and the induced operator on ternary residues
is the recurrence of {[}Tao22, Lemma 1.12{]} in density normalization.
Section 7.3 uses both as known. Its obstruction for finite weights, the
actual-integer error, and the two ancestor-mass criteria are supplied
here; no priority is claimed for the operator.

The companion papers {[}A{]}, {[}B{]}, {[}C{]}, and {[}D{]} treat
Juggler cycle bounds, finite-depth descent statistics, fate contagion,
and bounded-run \(3n-1\) cycles. We cite their results without reissuing
their proofs. The present main theorems require neither their large
computational floors nor the analytic estimates in {[}B{]}.

\subsection{1.2. Proof status}\label{proof-status}

This is the living preprint, version 0.8.1. It has not been
independently refereed; the availability section lists its Zenodo
versions. Mathematical priority for the signed adaptation and the
isolated obstruction results remains subject to specialist review.

Theorems 2.1, 3.2, 4.1, 4.4, 5.1, 6.1, and 7.3, Corollaries 4.2-4.3, and
Propositions 7.4-7.5 have compiled Lean statements; Appendix B
identifies their precise scope. Remark 5.5's larger exponent rests on
the same certificate and an exact integer comparison, not on a separate
Lean statement. Theorem 4.1 is unconditional in Lean: first-derivative
estimates, mixed-power cancellation, and the Fourier criterion prove its
simultaneous-box recurrence input. Theorem 4.4 also has a complete
quantitative Lean proof, including the finite derivative tests, every
cutoff Fourier mode, the half-open box estimate, both counting errors,
and the bounded witness. Appendix C gives an alternative written
argument using {[}AR24{]}; the formal proof independently derives
sufficient third- and fifth-derivative estimates and does not import the
exact external formula. The finite certificate is checked with exact
integers both in Lean and by an independent Python verifier. The
numerical search that found its weights is outside the proof. Kernel
checking, agreement between prose and formal statements, and independent
mathematical review are distinct checks.

\section{2. Words and the exact orbit
code}\label{words-and-the-exact-orbit-code}

For \(w=(b_0,\ldots,b_{L-1})\in\{0,1\}^L\), with \(0\) denoting E and
\(1\) denoting O, write \(o(w)=\sum_i b_i\) and \[
\rho(w)=\frac{3^{o(w)}}{2^L},\qquad
A(w)=\sum_{i=0}^{L-1}b_i2^i3^{\sum_{j>i}b_j}.
\tag{2.1}
\] Induction on the prescribed affine branches gives \[
2^L C_{\sigma,w}(x)=3^{o(w)}x+\sigma A(w).
\tag{2.2}
\] This describes an actual orbit only when all branch guards hold. For
an actual Juggler itinerary \(w\) starting at \(n\), the inequalities
\(J(n)^2\le n\) on E and \(J(n)^2\le n^3\) on O give \[
\bigl(J^L(n)\bigr)^{2^L}\le n^{3^{o(w)}}.
\tag{2.3}
\] Thus a contracting prefix forces descent for \(n>1\). For positive
\(3n+1\), the nonnegative correction in (2.2) instead prevents descent
along a prefix with \(\rho\ge1\). Shared combinatorics does not remove
this sign difference.

Let \(d_0<d_1<\cdots\) be the odd times of the Juggler orbit of \(n\),
finite or infinite, and define \[
H(n)=\sum_j\frac{2^{d_j}}{3^{j+1}}\quad\text{in }\mathbb Z_2.
\tag{2.4}
\] An empty sum is zero. The summands' valuations tend to infinity, so
the sum converges 2-adically. Both \(C_\sigma\) extend to
\(\mathbb Z_2\), using the residue modulo two to select the branch.
Formula (2.4) is a specialization of the inverse parity map of
{[}BL96{]}.

\textbf{Theorem 2.1 (exact signed orbit code).} For every positive
integer \(n\), \[
H(n)\equiv n\pmod2,\qquad
H(J(n))=g(H(n)),\qquad
-H(J(n))=C_{+1}(-H(n)).
\tag{2.5}
\] Two starts have equal codes if and only if all their finite Juggler
parity itineraries agree. On every finite source set, counting a
prescribed length-\(d\) itinerary is exactly counting its associated
code residue modulo \(2^d\).

\emph{Proof.} Deleting an even first letter gives \(H(n)=2H(J(n))\).
Deleting an odd first letter gives \(H(n)=1/3+(2/3)H(J(n))\). These
identities give parity and the minus-step identity. Negation gives the
plus-step identity. Iterate. The classical parity bijection identifies
length-\(d\) words with residues modulo \(2^d\). Equality of codes
therefore gives equality of every prefix; conversely, equality of all
residues gives equality in \(\mathbb Z_2\). The finite-set assertion is
equality of its membership conditions. \(\square\)

The formal construction takes compatible finite residues and their
2-adic limit. Its equality with (2.4) is proved by a telescoping sum
whose remainder after time \(k\) has 2-adic norm at most \(2^{-k}\). The
formal odd-time index uses the number of earlier odd times as \(j\);
finite and empty odd-time families are included. This assumes neither
termination nor 2-adic continuity of Juggler on its original integer
inputs.

\textbf{Example 2.2 (rational values and collisions).} From \(H(1)=1\),
following \(3,5,11,36,6,2,1\) backwards gives \[
\begin{array}{c|rrrrrrr}
n&3&5&11&36&6&2&1\\ \hline
H(n)&83/27&37/9&17/3&8&4&2&1.
\end{array}
\] Also \(H(4)=H(6)=4\). Thus the code neither takes every positive
integer to an ordinary integer nor is injective. Every terminating start
has a positive rational code whose reduced denominator divides a power
of three, by backwards application of \(q\mapsto2q\) and
\(q\mapsto(2q+1)/3\) from \(1\).

\textbf{Corollary 2.3 (the distribution question is unchanged).}
Frequencies \(2^{-d}\) for all length-\(d\) Juggler words, for every
fixed \(d\), are equivalent to \[
\lim_{N\to\infty}\frac1N
\#\{1\le n\le N:H(n)\equiv r\pmod{2^d}\}=2^{-d}
\] for every \(d\) and code residue \(r\). This equivalence proves
neither condition. The formal statement transfers limits using exact
equality of the finite-source counts.

\section{3. Actual periods and ordinary
integrality}\label{actual-periods-and-ordinary-integrality}

\textbf{Lemma 3.1 (order within a code fibre).} If \(H(n)=H(m)\) and
\(n\le m\), then \(J^k(n)\le J^k(m)\) for every \(k\ge0\).

\emph{Proof.} Equal codes imply equal parity at each time. Each branch
function is nondecreasing. Induction applies the same branch to the two
current values. \(\square\)

\textbf{Theorem 3.2 (period preservation and integrality).} Suppose
\(J^L(n)=n\) with \(L>0\), and let \(w\) be that actual itinerary. Put
\(D=3^{o(w)}-2^L\). Then \(D\ne0\), and:

\begin{enumerate}
\def\labelenumi{\arabic{enumi}.}
\tightlist
\item
  \(H(n)=A(w)/D\) in the 2-adic field.
\item
  For every \(d\ge0\), \(g^d(H(n))=H(n)\) if and only if \(J^d(n)=n\).
  The analogous statement holds for \(C_{+1}\) at \(-H(n)\).
\item
  The least periods of these three points are equal.
\item
  \(H(n)\) is an embedded ordinary integer exactly when \(D\mid A(w)\).
  In that case its ordinary signed Collatz orbit has the same return
  times as the Juggler orbit.
\end{enumerate}

\emph{Proof.} Iteration of (2.5) gives \[
3^{o(w)}H(n)=2^LH(J^L(n))+A(w).
\] The odd integer \(D\) is nonzero and invertible in \(\mathbb Z_2\).
Closing the orbit gives the formula.

One implication about return times is (2.5). Conversely, if the code
returns after \(d\) steps, \(J^d\) preserves its fibre and is
nondecreasing there by Lemma 3.1. If \(J^d(n)>n\), every subsequent
application stays strictly above \(n\); if \(J^d(n)<n\), every
subsequent application stays strictly below \(n\). Both contradict
\((J^d)^L(n)=n\). Hence \(J^d(n)=n\). Equality of return times gives
equality of least periods.

Finally \(A/D\) is an ordinary integer exactly when \(D\mid A\).
Injectivity of the rational embedding into the 2-adic field makes this
necessary as well as sufficient. The integer maps agree with their
2-adic extensions on embedded integers. \(\square\)

The periodicity assumption on the Juggler start is essential. A periodic
code alone has not been shown to imply a periodic start. For a
nontrivial Juggler cycle, \(D>0\) by strict power-envelope expansion
{[}A{]}. Its code is a positive rational \(g\)-periodic point, and its
negative is a negative \(C_{+1}\)-periodic point. No result here
establishes \(D\mid A(w)\) for such a cycle.

\section{4. Modular return does not force
integrality}\label{modular-return-does-not-force-integrality}

\textbf{Theorem 4.1 (actual modular returns with large code
denominator).} Fix \(a,b\ge1\) with \(3^a>2^{a+b}\). For every integer
\(M\ge1\) and lower bound \(B\), infinitely many odd \(n>B\) have actual
itinerary \(O^aE^b\), with \[
J^j(n)\ge n\ (0\le j\le a+b),\qquad J^{a+b}(n)>n,\qquad
J^{a+b}(n)\equiv n\equiv1\pmod{2M}.
\tag{4.1}
\] The rational periodic code associated with this word has denominator
\[
q_{a,b}=\frac{3^a-2^{a+b}}{\gcd(3^a-2^a,\,2^b-1)}.
\tag{4.2}
\] These denominators are unbounded at fixed \(b\). In particular \(M\)
may be any fixed power of \(q_{a,b}\).

\emph{Proof.} Put \(s=1+2Mt\), \(n=s^{2^{a-1}}\), and
\(\beta=3^a/2^{b+1}\). The first \(a-1\) odd steps are exact: \[
J^i(n)=s^{3^i2^{a-1-i}}\quad(0\le i<a).
\] The next value is \(y_0=\lfloor s^{3^a/2}\rfloor\). Whenever the
subsequent sources are even, their values are
\(y_j=\lfloor s^{2^{b-j}\beta}\rfloor\), \(0\le j\le b\). Nested
square-root floors are exact, since
\(\lfloor\sqrt{\lfloor x\rfloor}\rfloor=\lfloor\sqrt x\rfloor\).

As \(t\) varies, the vector \[
\left(\frac{s^\beta}{2M},\frac{s^{2\beta}}2,
\frac{s^{4\beta}}2,\ldots,\frac{s^{2^b\beta}}2\right)\pmod1
\tag{4.3}
\] is uniformly distributed. Every nonzero integer linear combination
has a nonintegral leading exponent: all exponents have odd numerator and
denominator divisible by two. It is a Hardy-field function of polynomial
growth, and its distance from any rational polynomial, divided by
\(\log t\), tends to infinity. Boshernitzan's criterion {[}Bo94, Theorem
1.3{]}, also stated in {[}R26, Theorem 1.1{]}, applies. Weyl's criterion
gives joint distribution.

Place the first coordinate's fractional part in \([1/(2M),2/(2M))\), and
the others in \([0,1/2)\). The box has volume \(1/(2^{b+1}M)\). It
enforces \(y_b\equiv1\pmod{2M}\) and even \(y_0,\ldots,y_{b-1}\), so the
word is actual. Since \(\beta>2^{a-1}\), its endpoint exceeds \(n\) for
sufficiently large \(s\). The odd portion increases and the even portion
decreases to that endpoint, proving (4.1).

For \(O^aE^b\), (2.1) gives \(A=3^a-2^a\) and \(D=3^a-2^{a+b}\). The
oddness of \(A\) gives \(\gcd(A,D)=\gcd(A,2^b-1)\), proving (4.2). The
quotient is at least \((3^a-2^{a+b})/(2^b-1)\), which tends to infinity
with \(a\) at fixed \(b\). \(\square\)

The formal construction proves the exact root-floor identities, all
branch guards, and every conclusion of (4.1) for each successful
parameter with \(s\ge2^{2^{b+1}}\). It also proves (4.2) and, for every
\(Q\), the explicit sufficient bound \[
a\ge2\bigl(2^b+(Q+1)(2^b-1)\bigr)\quad\Longrightarrow\quad q_{a,b}>Q.
\] The formal analytic proof establishes BoxRecurrence: for every \(T\)
there is a parameter \(t\ge T\) in the simultaneous box above. It uses
an elementary derivative argument in place of the Hardy-field criterion.
For a phase with leading term \(cx^p\), \(c\ne0\), \(0<p<1\), the
first-derivative estimate gives cancellation. A mean-value point in each
fixed shifted interval proves that the actual difference lowers the
leading exponent by one; induction and van der Corput's inequality
handle higher noninteger powers. The same argument controls lower-order
terms and the progression \(s=1+2Mt\). Uniform density of Fourier
monomials on the torus, followed by a nonnegative continuous function
supported in the interior of the target box, supplies arbitrarily late
visits. Thus PaperERecurrence.theorem41 has no recurrence or
exponential-sum hypothesis.

The theorem concerns growing prefixes and their associated
\emph{periodic-word} codes. It does not identify that code with
\(H(n)\), which depends on the entire subsequent itinerary. It gives no
common start at all precisions, no actual nontrivial cycle, and no
denominator bound for true Juggler cycles. Finite modular return is not
exact orbit closure.

\textbf{Corollary 4.2 (counting the constructed returns).} Fix \(a,b,M\)
as in Theorem 4.1 and put \(d=2^{a-1}\). Let \(\mathcal T\) be the set
of nonnegative parameters \(t\) for which (4.3) lies in the half-open
box used above and \(s=1+2Mt\ge2^{2^{b+1}}\). Define \[
R_{a,b,M}(X)=\#\{t\in\mathcal T:(1+2Mt)^d\le X\}.
\] Then, as \(X\to\infty\), \[
R_{a,b,M}(X)\sim\frac{X^{1/d}}{2^{b+2}M^2}.
\tag{4.4}
\] For every fixed \(\varepsilon>0\), every sufficiently large interval
\((X,(1+\varepsilon)X]\) contains a start from this family and hence an
actual modular return satisfying (4.1).

\emph{Proof.} Uniform distribution and the box volume give \[
\frac{\#(\mathcal T\cap\{0,\ldots,N-1\})}{N}
\longrightarrow\delta:=\frac1{2^{b+1}M}.
\] The threshold excludes only finitely many parameters. For \(X\ge1\),
the number of possible parameter positions is exactly \[
N(X)=\left\lfloor\frac{X^{1/d}-1}{2M}\right\rfloor+1
\sim\frac{X^{1/d}}{2M}.
\] Distinct parameters give distinct starts. Substituting \(N(X)\) into
the density limit yields (4.4). The difference
\(R_{a,b,M}((1+\varepsilon)X)-R_{a,b,M}(X)\), divided by \(X^{1/d}\),
tends to \(\bigl((1+\varepsilon)^{1/d}-1\bigr)/(2^{b+2}M^2)>0\). It is
therefore positive for every sufficiently large \(X\). \(\square\)

This counts the explicit perfect-power family, not all starts satisfying
(4.1). In particular, it is not a positive-density assertion among all
positive integers.

\textbf{Corollary 4.3 (prescribed residues in the even run).} Fix the
same \(a,b,M\), and choose residues \(0\le r_j<2M\), \(0\le j\le b\),
with \(r_j\) even for \(j<b\) and \(r_b=1\). Infinitely many starts
\(n=(1+2Mt)^{2^{a-1}}\) satisfy (4.1) and \[
J^{a+j}(n)\equiv r_j\pmod{2M}\qquad(0\le j\le b).
\tag{4.5}
\] More precisely, the constructing parameters beyond the same
threshold, defined by \[
\left\{\frac{s^{3^a/2^{j+1}}}{2M}\right\}
\in\left[\frac{r_j}{2M},\frac{r_j+1}{2M}\right)
\quad(0\le j\le b),
\tag{4.6}
\] have natural density \((2M)^{-(b+1)}\) in the parameter \(t\).

\emph{Proof.} The distinct positive noninteger exponents are unchanged;
scaling every coordinate by \(1/(2M)\) preserves the nonzero-mode
argument. Joint uniform distribution gives the product of the \(b+1\)
interval lengths, hence the asserted density. For every nonnegative real
\(x\), the fractional-part condition
\(\{x/(2M)\}\in[r/(2M),(r+1)/(2M))\) is equivalent to
\(\lfloor x\rfloor\equiv r\pmod{2M}\). Thus (4.6) prescribes the exact
root values \(y_j\). The even residues supply every branch guard; the
final residue gives the exit congruence. The threshold ensures expansion
and all the inequalities in (4.1). Removing its finite initial segment
preserves density, and the strictly increasing parameter-to-start map
gives infinitude above every bound. \(\square\)

The half-open endpoints, including residue zero, are included in the
formal frequency proofs. Fourier convergence yields weak convergence of
the empirical measures to Haar measure. The boundaries of these boxes
lie in finitely many coordinate endpoint sets of Haar measure zero, so
their frequencies equal their volumes. This strengthens the open-box
recurrence argument without assuming a new analytic estimate. Both
corollaries fix \(a,b,M\), the residues, and any \(\varepsilon>0\)
before taking limits. They provide neither a bound on the first witness
nor estimates uniform in growing parameters, and do not concatenate the
finite prefixes into an infinite orbit.

\subsection{4.1. Effective returns for the word
OOE}\label{effective-returns-for-the-word-ooe}

The preceding results fix the modulus before taking a limit. The next
result retains its dependence explicitly, for one fixed word. For
integers \(M,T\ge1\), put \(s_t=1+2Mt\) and \[
A_M(T)=\#\{0\le t<T:s_t\ge16,\quad
 \lfloor s_t^{9/2}\rfloor\equiv0\pmod2,\quad
 \lfloor s_t^{9/4}\rfloor\equiv1\pmod{2M}\}.
\]

\textbf{Theorem 4.4 (effective OOE returns).} For every \(M,T\ge1\), \[
\begin{split}
\left|A_M(T)-\frac{T}{4M}\right|
&\le\left[5+128M^{1/4}\left(3+\frac{\log T}{16}\right)^2\right]T^{63/64}+8\\
&\le2^{14}M^{1/4}T^{127/128}.
\end{split}
\tag{4.7}
\] In particular, some \(0\le t<2^{2176}M^{160}\) gives a start
\(n=s_t^2<2^{4354}M^{322}\) whose actual three-step itinerary is OOE,
whose intermediate states are at least \(n\), and whose exit is greater
than \(n\) and congruent to \(n\equiv1\pmod{2M}\). At
\(T=2^{2176}M^{160}\), there are at least \(T/(8M)\) such parameters.

\emph{Proof.} Appendix C proves the uniform Fourier estimates, the
finite half-open box bound, their dyadic assembly, and the witness
extraction. The box is \([0,1/2)\times[1/(2M),1/M)\); its area is
\(1/(4M)\), including when \(M=1\). The exact construction of Theorem
4.1 with \(a=2,b=1\) converts these parameter conditions into the
asserted orbit conditions. \(\square\)

The constants are theoretical bounds, not practical search budgets. This
OOE word has periodic-word denominator one; Theorem 4.4 does not make
the large-denominator family of Theorem 4.1 effective. Both the counting
estimate and the actual bounded return are kernel-checked. Independent
mathematical review remains outstanding.

\section{5. A signed ancestor-count
theorem}\label{a-signed-ancestor-count-theorem}

Define \[
\pi^-_a(X)=\#\{1\le n\le X:\exists j\ge0,\ g^j(n)=a\}.
\] Let \(N(a,X)\) additionally require \(g^i(n)\le X\) for every
\(0\le i\le j\). Then \(N(a,X)\le\pi^-_a(X)\); both counts are
nondecreasing in \(X\). They count starts, not paths.

\textbf{Theorem 5.1 (positive \(3n-1\) ancestors).} For every positive
integer \(a\) with \(3\nmid a\), there is \(X_0(a)\) such that, for
every integer \(X\ge X_0(a)\), \[
X^{21}\le N(a,X)^{25}\le\bigl(\pi^-_a(X)\bigr)^{25}.
\tag{5.1}
\] In particular \(\pi^-_a(X)\ge X^{21/25}\).

\subsection{5.1. Why height changes with the
sign}\label{why-height-changes-with-the-sign}

A fertile \(g\)-target has \(a\equiv1\pmod3\) and odd predecessor
\(b=(2a+1)/3>2a/3\). For \(C_{+1}\), the odd predecessor at
\(a\equiv2\pmod3\) is \((2a-1)/3<2a/3\). The homogeneous child budget
\((X/a)(3/2)b\) therefore exceeds \(X\) on the minus side. For \(a=19\),
\(b=13\), \(X=103\), the nominal budget is \(4017/38>105\), admitting
\(104,52,26,13\) even though its start already exceeds the actual
cutoff.

\subsection{5.2. Valid grid recurrences}\label{valid-grid-recurrences}

Let \(r_0,\ldots,r_{49}\) be the integer table in Appendix A and set \[
C(t)=2^{\lfloor t/50\rfloor}r_{t\bmod50},\qquad
F_a(t)=N(a,C(t)a).
\] Thus \(C(t+100)=4C(t)\) and \(C(50j)=10000\,2^j\). Integer checks on
the fifty residues of \(t\) give \[
8193C(t+29)\le12288C(t),\qquad
16386C(t)\le12288C(t+21).
\tag{5.2}
\] If \(a\ge4096\) and \(3b=2a+1\), then \(12288b\le8193a\). Combining
these inequalities puts every selected child cutoff below its parent
cutoff.

\textbf{Lemma 5.2 (actual inverse trees).} If \(a\ge4096\),
\(a\equiv1\pmod3\), and \(a\) is nonperiodic, then \[
\begin{aligned}
F_a(t+100)&\ge F_{4a}(t),\\
F_a(t+100)&\ge F_{4a}(t)+F_b(t+129),\\
F_a(t+100)&\ge F_{4a}(t)+F_{2b}(t+79).
\end{aligned}
\tag{5.3}
\] The first inequality needs neither nonperiodicity nor the threshold.

\emph{Proof.} The child paths are \(4a,2a,a\), \(b,a\), and \(2b,b,a\).
The cutoff comparisons permit extension to \(a\). The \(4a\) subtree is
disjoint from each alternative odd subtree. Otherwise a common start
reaches the nonperiodic root at two different times, or the last steps
force odd \(b\) to equal even \(2a\). A nonperiodic root has a unique
hitting time: two times would give it a positive return. Thus the
indicated cardinalities add. \(\square\)

Use the second recurrence at \(a\equiv1\pmod9\), the third at
\(a\equiv7\pmod9\), and only the first at \(a\equiv4\pmod9\). These
choices keep children in residue \(1\pmod3\). The two odd alternatives
are never added together.

Set \(\mathcal M(t,a)=10t+\lfloor\log_2(a^{498})\rfloor\). The exact
inequality \[
2^{291}8193^{498}<12288^{498}
\tag{5.4}
\] forces the odd child's root-size term to fall by at least 291.
Doubling or quadrupling adds exactly 498 or 996. Therefore \[
\begin{aligned}
\mathcal M(t,4a)+4&\le\mathcal M(t+100,a),\\
\mathcal M(t+79,2b)+3&\le\mathcal M(t+100,a),\\
\mathcal M(t+129,b)+1&\le\mathcal M(t+100,a).
\end{aligned}
\tag{5.5}
\] Even the advanced-index call decreases this natural-number measure.

\subsection{5.3. A closed domain for arbitrary
targets}\label{a-closed-domain-for-arbitrary-targets}

\textbf{Lemma 5.3 (root domain).} Every positive target \(a\) prime to
three has a nonperiodic ancestor \(r>2^{19}\) with \(r\equiv1\pmod3\).
All fertile ancestors of \(r\) are at least 4096 and form a domain
closed under the selected productions in (5.3). For all sufficiently
large \(X\), \(N(r,X)\le N(a,X)\).

\emph{Proof.} At most three doublings reach a value \(z\equiv1\) or
\(7\pmod9\). Its two predecessors \(2z\) and \((2z+1)/3\) are distinct,
positive, and prime to three. They cannot both be periodic, since a
deterministic map is injective on its periodic points. Select a
nonperiodic predecessor, double once if necessary to reach residue 1
modulo three, then multiply by \(2^{20}\). This gives \(r\) as required.
Ancestors of nonperiodic points are nonperiodic.

The finite barrier used here is \[
0\le u<4096\quad\Longrightarrow\quad
g^k(u)<2^{19}\quad(k\ge0),
\tag{5.6}
\] where \(g(0)=0\). A finite descent certificate and strong induction
prove it: each start reaches a smaller value or the explicit
forward-invariant set of zero and the fifteen members of the three known
positive cycles. Appendix A specifies an independent finite check. No
exhaustiveness of the known cycles is assumed. An ancestor below 4096
could not reach \(r>2^{19}\), by (5.6). The residue calculations after
(5.3) give closure. Above the maximum on the fixed path from \(r\) to
\(a\), every capped path to \(r\) extends to a capped path to \(a\).
\(\square\)

A threshold alone is insufficient: 4096 has odd predecessor 2731.
Restricting to ancestors of \(r\) repairs that boundary.

\subsection{5.4. Certificate and growth
induction}\label{certificate-and-growth-induction}

Take \(p=5059\), \(q=5000\), \(C_{\max}=10^{12}\), \(\mu=p/q\). The
certificate has a positive integer weight \(c_m\) for each of the
\(3^{11}=177147\) fertile classes modulo \(3^{12}\). Its minimum is
7307142888 and its maximum is \(C_{\max}\). For a child class modulo
\(3^{11}\), let \(\bar c\) be the minimum over its three lifts modulo
\(3^{12}\). Every row satisfies the exact integer inequality \[
\begin{aligned}
c_m p^{100}&\le c_{4m}q^{100}
&& (m\equiv4\pmod9),\\
c_m p^{100}&\le c_{4m}q^{100}
 +\bar c_{(4m+2)/3}p^{79}q^{21}
&& (m\equiv7\pmod9),\\
c_m p^{100}q^{29}&\le c_{4m}q^{129}
 +\bar c_{(2m+1)/3}p^{129}
&& (m\equiv1\pmod9).
\end{aligned}
\tag{5.7}
\] Indices are reduced at the indicated moduli. A minimum ensures
validity for the actual child's lift, irrespective of root size.

\textbf{Lemma 5.4 (growth).} Throughout the domain of Lemma 5.3, \[
F_a(t)\ge\frac{c_a}{C_{\max}}\mu^{t-100}\quad(t\ge0).
\tag{5.8}
\]

\emph{Proof.} Induct on \(\mathcal M(t,a)\). For \(t<100\) the root
itself is counted and the right side is at most one. For \(t=s+100\),
use (5.3), the decreases (5.5), and the inductive bounds. The three
cases reduce exactly to (5.7) after division by the appropriate positive
powers of \(p,q\). The minimum weight does not exceed the actual child's
weight. Equivalently the induction uses only integers:
\(c_ap^tq^{100}\le F_a(t)C_{\max}q^tp^{100}\). \(\square\)

\emph{Proof of Theorem 5.1.} Apply (5.8) at \(r\) and \(t=50j\). For all
sufficiently large \(j\), Lemma 5.3 gives \[
N(a,10000r\,2^j)\ge K\mu^{50j}
\] with fixed \(K>0\). Exact integer comparison gives \[
2^{21}5000^{1250}<5059^{1250},
\tag{5.9}
\] so \(\mu^{50}>2^{21/25}\). For \(10000r\,2^j\le X<10000r\,2^{j+1}\),
monotonicity gives the same lower bound at \(X\). The strict exponential
advantage eventually absorbs \(K\) and the fixed interpolation factor.
Raise to the twenty-fifth power to obtain (5.1). \(\square\)

This is a lower bound on ancestors, not positive natural density or a
conclusion about every orbit. No classification of all positive \(3n-1\)
cycles enters its proof.

\textbf{Remark 5.5 (the certificate's full exponent).} The rate
\(\mu=5059/5000\) gives more than (5.9) records:
\(50\log_2\mu=0.84620\ldots\), and exact integer comparison gives \[
2^{423}5000^{25000}<5059^{25000}<2^{424}5000^{25000}.
\tag{5.10}
\] The proof of Theorem 5.1 then yields \(X^{423}\le N(a,X)^{500}\) for
all sufficiently large \(X\), so \(\pi^-_a(X)\ge X^{423/500}\). The
second inequality in (5.10) shows that \(423/500\) is the best exponent
of this form from this certificate. The Lean statement keeps \(21/25\).

\section{6. A common ceiling for the fixed
grid}\label{a-common-ceiling-for-the-fixed-grid}

Keep shifts \(100,79,129\) fixed, but allow any finite residue level
\(\ell\ge2\). Put \(M=3^{\ell-2}\) and take positive weights
\(c_0,\ldots,c_{3M-1}\). For minus use representatives \(3i+1\); for
plus use \(3i+2\). The normalized minus rows are (5.7) with rate
\(\mu>0\); their coefficients are \(\mu^{-100},\mu^{-21},\mu^{29}\). For
plus the odd children are \((4m-2)/3\) at \(m\equiv2\pmod9\) and
\((2m-1)/3\) at \(m\equiv8\pmod9\); the sterile parent is
\(m\equiv5\pmod9\). The same coefficients define its rows.

\textbf{Theorem 6.1 (fixed-grid ceiling for both signs).} If the rows
hold for either sign at any finite level, and \(0<\mu\) with
\(\mu^{50}\le2\), then \[
\mu<5069/5000,\qquad \mu^{5000}<2^{99}.
\tag{6.1}
\] Thus \(\gamma=50\log_2\mu<99/100\). In particular the harmonic rate
\(\mu^{50}=2\) is impossible.

\emph{Proof.} Put \[
S=\sum_{i=0}^{3M-1}c_i,\qquad
B=\sum_{j=0}^{M-1}\min(c_j,c_{j+M},c_{j+2M}).
\] Then \(3B\le S\). The fourfold-parent index maps are \(i\mapsto4i+1\)
for minus and \(i\mapsto4i+2\) for plus, modulo \(3M\). The odd-child
maps modulo \(M\) are \(j\mapsto2j,\ j\mapsto4j+3\) for minus, and
\(j\mapsto2j+1,\ j\mapsto4j\) for plus. All are permutations. Summing
the rows gives \(S\le\mu^{-100}S+(\mu^{29}+\mu^{-21})B\), hence \[
1\le F(\mu):=\mu^{-100}+\frac{\mu^{29}+\mu^{-21}}3.
\tag{6.2}
\] The function \(F\) is convex on the positive reals, as each power has
nonnegative second derivative. Exact arithmetic gives \[
\begin{gathered}
F(5069/5000)<1,\quad F(507/500)<1,\\
(5069/5000)^{50}<2<(507/500)^{50},\\
(5069/5000)^{5000}<2^{99}.
\end{gathered}
\tag{6.3}
\] Convexity makes \(F<1\) between the rational endpoints. Every
\(\mu\ge5069/5000\) with \(\mu^{50}\le2\) lies there, contradicting
(6.2). The last comparison proves (6.1).

At the harmonic rate, clearing (6.2) gives
\(3\mu^{100}\le3+\mu^{129}+\mu^{79}\). If \(\mu^{50}=2\), it forces
\(\mu^{79}\ge3\) and then \(2^{79}\ge3^{50}\), contrary to
\(2^{79}<3^{50}\). \(\square\)

This obstruction reflects the fixed rounding slack \(79/50<\log_2 3\).
It applies at all finite table sizes, but not to every finer grid,
different production system, or direct reciprocal-mass argument. It is
not an upper bound for actual ancestor counts.

\section{7. Two other limits of
transfer}\label{two-other-limits-of-transfer}

\subsection{7.1. Counts and reciprocal
mass}\label{counts-and-reciprocal-mass}

\textbf{Proposition 7.1 (finite backward mass).} For either shortcut
sign, the set \(\mathcal R=\{3\cdot2^k:k\ge0\}\) is infinite and
backward-closed on the positive integers, with
\(\sum_{n\in\mathcal R}1/n=2/3\).

\emph{Proof.} A multiple of three has no odd predecessor:
\(3u+\sigma=2m\) would imply \(\sigma\equiv0\pmod3\). Its only
predecessor is its double. Sum the geometric series
\(\frac13\sum_{k\ge0}2^{-k}\). \(\square\)

The set is not forward-closed: \(C_{+1}(3)=5\), \(g(3)=4\). It refutes a
statement for all backward-closed sets, not one restricted to complete
fate classes. Juggler's even inverse block consists of the even integers
in \([m^2,(m+1)^2)\). Its reciprocal mass has the uniform lower bound
\(m/(m+1)^2\) used in {[}C{]}. The code preserves neither these ordinary
heights nor their reciprocal weights.

A lower bound \(N(X)\gg X^\kappa\) with \(0<\kappa<1\) alone cannot
force divergent reciprocal mass. For example the distinct integers
\(\lfloor j^{1/\kappa}\rfloor\), \(j\ge1\), have count at least
\(\lfloor X^\kappa\rfloor\) and convergent reciprocal sum. Theorem 5.1
therefore supplies no missing harmonic estimate for a fate-class
argument.

\subsection{7.2. Stopping a word family}\label{stopping-a-word-family}

At fixed depth, fair words satisfy \[
\sum_{|w|=d}2^{-d}=1,\qquad
\sum_{|w|=d}2^{-d}\rho(w)
=\left(\frac12\frac12+\frac12\frac32\right)^d=1.
\tag{7.1}
\] Both identities persist for bounded complete prefix trees. The formal
statement represents these as finite binary trees with both children at
every internal node, and sums over their leaf words. Completeness alone
does not suffice at unbounded depth.

\textbf{Proposition 7.2 (moment loss at unbounded stopping).} Let
\(\mathcal C\) consist of the minimal words whose multiplier first
becomes less than one. It is prefix-free, and \[
\sum_{w\in\mathcal C}2^{-|w|}=1,\qquad
\sum_{w\in\mathcal C}2^{-|w|}\rho(w)\le3/4.
\tag{7.2}
\]

\emph{Proof.} The fair log-multiplier walk has mean increment
\(\frac12\log(3/4)<0\). The strong law makes its first strictly negative
time finite almost surely. Its minimal stopping prefixes partition this
probability-one event. Every stopped multiplier is less than one; E
alone has mass \(1/2\) and multiplier \(1/2\), losing \(1/4\) relative
to fair mass. The other terms cannot restore it. \(\square\)

The Lean proof instead uses the established survivor-count decay for
completeness. It then regroups the nonnegative sums by word length to
prove (7.2) directly on the type of minimal certificate words. The
probabilistic proof explains the missing moment without invoking an
optional-stopping equality.

\subsection{7.3. Reciprocal mass on signed
fibres}\label{reciprocal-mass-on-signed-fibres}

For \(s\in\{1,-1\}\) let \(S_s(n)=(3n+s)/2^{v_2(3n+s)}\) on the positive
odd integers, the odd-to-odd acceleration of \(C_s\). A positive odd
target \(m\) has no odd predecessor if \(3\mid m\). Otherwise its odd
predecessors are \(n_k=(2^km-s)/3\) for the \(k\ge1\) with
\(2^km\equiv s\pmod3\); they form one ray under \(n\mapsto4n+s\), and
each run of \(3^r\) consecutive members meets every residue class modulo
\(3^r\) once {[}Z08{]}. For a real weight \(h\) on the residues modulo
\(3^r\), vanishing on multiples of three, put \[
(L_sh)(a)=\sum_{\substack{k\ge1\\2^ka\equiv s\ (3)}}\frac3{2^k}\,
h\Bigl(\frac{2^ka-s}3\Bigr),
\tag{7.3}
\] a function of \(a\) modulo \(3^{r+1}\). This is the recurrence of
{[}Tao22, Lemma 1.12{]}. For \(0\le h\le H\) and every positive odd
\(m\), the actual weighted mass \(W_h(m)=\sum h(n)/n\) over the odd
\(n\) with \(S_s(n)=m\) satisfies \[
\bigl|mW_h(m)-(L_sh)(m)\bigr|\le\frac H{m-1/2}.
\tag{7.4}
\] Multiplying a predecessor's term by \(m\) gives
\((3/2^k)h(n_k)/(1-s/(2^km))\); subtracting its homogeneous part and
summing \(3H/(4^km)\) over \(k\ge1\) gives (7.4).

\textbf{Theorem 7.3 (no finite ternary weight reproduces).} Let
\(r\ge1\), let \(h\) be a weight modulo \(3^r\) that vanishes on
multiples of three and is positive on the units, and let \(d\ge1\). For
either sign, some unit residue \(a\) modulo \(3^{r+d}\) satisfies \[
(L_s^dh)(a)<h(a\bmod3^r).
\tag{7.5}
\] Consequently the positive odd units \(m\) with \(mW_h(m)<h(m)\) have
divergent reciprocal sum, and every set of finite reciprocal mass misses
one of them. More generally, if \(h\ge0\) vanishes on multiples of three
and \(h(-s)>0\), then for every \(D\ge0\) one unit residue modulo
\(3^{r+D+1}\) defeats every inverse-tree policy that expands the first
generation and stops each branch by depth \(D+1\): its homogeneous
payoff is below \(h\) at every positive integer in that residue.

\emph{Proof of (7.5).} For each \(k\), \(b\mapsto2^{-k}(3b+s)\) is a
bijection between the child classes and the admissible parent classes,
so summing (7.3) over \(a\) gives \(\sum_aL_sh(a)=3\sum_bh(b)\), and
iteration gives the factor \(3^d\). The same total is the sum of
\(h(a\bmod3^r)\) over the classes modulo \(3^{r+d}\). If no unit row
were deficient, every row would be an equality, since nonunit rows are
zero on both sides. At \(a\equiv-s\) the branch \(k=1\) returns \(-s\)
at the next lower modulus with coefficient \(3/2\), so
\((L_s^dh)(-s)\ge(3/2)^dh(-s)>h(-s)\), a contradiction. \(\square\)

For \(d=1\), a deficiency \(\delta>0\) and (7.4) make every odd
\(m\ge1/2+2H/\delta\) in the deficient class satisfy
\(mW_h(m)\le h(m)-\delta/2\), an arithmetic progression of divergent
reciprocal sum. The policy statement applies the summation to the
envelope \(U_0=h\), \(U_{j+1}=\max(h,L_sU_j)\), which dominates every
policy of the stated depth and is attained by one; its proof is formal
(Appendix B). No finite residue weight, table size, grouping depth, or
bounded stopping rule therefore gives uniform reproduction of reciprocal
mass. This is the Collatz counterpart of the uniform production that
Paper {[}C{]} obtains on Juggler fibres. It does not bound the
intersection of the deficient targets with any fate class.

\textbf{Proposition 7.4 (ancestor-mass criteria).} Let \(a\) be a
positive odd integer that is not periodic under \(S_s\), and put
\(K_d(a)=(L_s^d\mathbf 1_{\mathrm{units}})(a)\), where
\(\mathbf 1_{\mathrm{units}}\) is the indicator of the units modulo
three. If \(\sum_{d\ge0}K_d(a)=\infty\), then \(\sum1/n\) over the
positive odd ancestors \(n\) of \(a\) prime to three diverges, for
either sign.

\emph{Proof.} A depth-\(d\) inverse path with exponents
\(k_1,\ldots,k_d\) has a unique endpoint, and its coefficient in
\(K_d(a)\) is \(3^d/2^{k_1+\cdots+k_d}\). For \(s=1\) each inverse step
subtracts \(1/3\) before dividing by three, so the endpoint is at most
\(a2^{\sum k_i}/3^d\) and generation \(d\) has reciprocal mass at least
\(K_d(a)/a\). For \(s=-1\) the endpoint exceeds that value by the factor
\(\prod_j(1-1/(3x_j))^{-1}\) over the path's odd states \(x_j\).
Distinct states of a nonrepeating finite path have reciprocal sum below
one absolute constant \(B\), by the exact binary-residue moment of the
shortcut map and an equal-time packing, so generation \(d\) has mass at
least \(K_d(a)e^{-B}/a\). A nonperiodic target is reached at one time
only, so the generations are disjoint and their masses add. \(\square\)

The divergence premise is open at every target; the proposition is a
reduction, not a lower bound. Every nonempty full fate class of
\(C_{+1}\) contains such a target: a forward iterate is an odd unit
whose fibre has infinitely many unit members, at most one of them
periodic. Pairing the two signs does not help at low depth.

\textbf{Proposition 7.5 (no cross-sign compensation at depth two).} Read
at every positive integer, the formula for \(S_s\) gives
\(S_-(2n+1)=S_+(n)\) and \(S_+(2n-1)=S_-(n)\). At every
\(a\equiv4\pmod{27}\), \[
K_2^+(a)=\frac{12076}{29127},\qquad K_2^-(a)=\frac{7988}{29127},\qquad
K_2^+(a)+K_2^-(a)<1 .
\] If \(M_2^s(a)\) is the actual reciprocal mass of the positive odd
sources prime to three reaching \(a\) in exactly two returns of \(S_s\),
then \(a(M_2^+(a)+M_2^-(a))<3/4\) for every \(a=31+54t\), \(t\ge0\).

The identities are direct. The coefficients are exact rational sums over
the periodic exponent classes, and the actual masses follow from a
two-step form of (7.4). Iterating the paired operator introduces
mixed-sign paths, and no statement over all depths is made.

\section{8. What remains open}\label{what-remains-open}

Integer transport on an actual Juggler cycle requires \(D\mid A(w)\), or
a useful restriction on its reduced denominator. Distribution of Juggler
words requires estimates of integer-source multiplicities in code
cylinders. Transport of fate contagion requires control of ordinary
reciprocal mass. Existence of the exact code supplies none of these
additional estimates. On the Collatz side, Section 7.3 rules out
finite-weight reproduction and leaves one arithmetic question: whether
\(\sum_dK_d(a)\) diverges at a nonperiodic target in each fate class.
Averaged residue mixing controls densities, not these prescribed integer
targets.

The signed ancestor theorem leaves universal termination and the
classification of all cycles open. Its grid ceiling is a restriction of
the stated certificate family. The finite-mass ray does not settle
harmonic divergence for a full fate class.

Future versions may add results with their hypotheses, proof status,
prior-art comparison, and reproducibility record. The companion papers
retain their own claims and histories.

\section{Acknowledgments and use of
AI}\label{acknowledgments-and-use-of-ai}

The author thanks the authors of the cited work and the Lean and Mathlib
communities for the mathematical and formal infrastructure. M. Sharpe's
grid implementation is specifically acknowledged in Section 1.1 and the
attributed source modules.

AI assistance was used in mathematical exploration, drafting,
programming, and formalization. The author is responsible for the
manuscript and its claims. Kernel-checked results are identified
individually; AI assistance and local checking do not replace
independent mathematical review.

\section{Availability and versioning}\label{availability-and-versioning}

The project is maintained at
\href{https://github.com/sneakyweasel/btlab}{github.com/sneakyweasel/btlab}.
The canonical source is
\texttt{\nolinkurl{docs/theory/juggler_signed_collatz_note.md}}. Build
with the command given in Appendix A; the update and validation guide is
\texttt{\nolinkurl{docs/theory/PAPER_E_BUILD.md}}. The release manifest
records SHA-256 hashes of the manuscript, mathematical dependencies,
build tools, and outputs. Version 0.8.1 adds these DOIs and the common
build commands to the text of 0.8.0. The Zenodo versions are 0.8.0 of 25
September 2026
(\href{https://doi.org/10.5281/zenodo.22954746}{doi:10.5281/zenodo.22954746})
and 0.7.1 of 22 September 2026
(\href{https://doi.org/10.5281/zenodo.22905650}{doi:10.5281/zenodo.22905650});
the concept DOI
\href{https://doi.org/10.5281/zenodo.22905649}{10.5281/zenodo.22905649}
resolves to the latest version.

The manuscript uses CC BY 4.0. Software retains its repository license
and third-party notices; reused grid material retains M. Sharpe's MIT
notice. The local deposit kit contains the PDF and a
source-and-certificate archive, so finite proof data need not be
recovered from a moving repository branch.

\section{Appendix A. Finite data and
reproduction}\label{appendix-a.-finite-data-and-reproduction}

The cap table is indexed from zero:

\begin{verbatim}
10000 10140 10281 10425 10570 10718 10867 11019 11173 11329
11487 11647 11810 11975 12142 12311 12483 12658 12834 13013
13195 13379 13566 13755 13947 14142 14340 14540 14743 14948
15157 15369 15583 15801 16021 16245 16472 16702 16935 17171
17411 17654 17901 18150 18404 18661 18921 19185 19453 19725
\end{verbatim}

The full table is stored in
\texttt{\nolinkurl{data/research/juggler/negative_preimage_density/grid_k12_certificate.json}}.
Its ordering is \(m=3i+1\), \(0\le i<177147\). For each row compute
\(4m\bmod3^{12}\), and, when needed, the odd child modulo \(3^{11}\).
Read the minimum over its three lifts, then check (5.7) with integers.
Recorded floating-point diagnostics are ignored.

The weights were found by damped iteration of the positive production
operator from all ones, normalizing the maximum and rounding at scale
\(10^{12}\). The first successful table occurred after 62 iterations.
Any table passing the exact checks would suffice; its discovery is
outside the proof.

For an independent check of (5.6), start from all integers below 4096
and repeatedly add successors until no new state occurs. The closure of
the positive starts has 6417 states and maximum 417718, below
\(2^{19}=524288\); zero is fixed and can be added separately. The
verifier checks closure and the strict barrier.

\begin{verbatim}
python tools/check_paper_e.py
python tools/generate_signed_grid_certificate.py --check
python tools/build_paper.py E
python tools/build_paper.py E --check
cd formal
lake build Problems.JugglerCollatzPaper
lake env lean AxiomCheckJugglerCollatzPaper.lean
\end{verbatim}

The first command checks the integer certificate, grid, barrier,
rational examples, modular-return witnesses, and exact comparisons
(5.4), (5.9), and (6.3), together with the manuscript's numbered
references and declaration inventory. These finite checks do not prove
Theorem 4.1's infinite equidistribution assertion. The second command
verifies the five generated Lean certificate modules against the data.

\clearpage

\section{Appendix B. Formalization and review
boundary}\label{appendix-b.-formalization-and-review-boundary}

The paper barrel is Problems.JugglerCollatzPaper. It imports the
existing proofs without changing their hypotheses.

\begingroup\small

\begin{longtable}[]{@{}
  >{\raggedright\arraybackslash}p{(\linewidth - 2\tabcolsep) * \real{0.3600}}
  >{\raggedright\arraybackslash}p{(\linewidth - 2\tabcolsep) * \real{0.6400}}@{}}
\toprule\noalign{}
\begin{minipage}[b]{\linewidth}\raggedright
Result
\end{minipage} & \begin{minipage}[b]{\linewidth}\raggedright
Module and declaration
\end{minipage} \\
\midrule\noalign{}
\endhead
\bottomrule\noalign{}
\endfoot
\bottomrule\noalign{}
\endlastfoot
Theorem 2.1 & \texttt{\nolinkurl{CollatzPadic.orbit_bridge}} \\
Series (2.4) &
\texttt{\nolinkurl{PaperECompletion.code_hasSum_odd_times}} \\
Codes and cylinders & \texttt{\nolinkurl{CollatzPadic.code_eq_iff}},
\texttt{\nolinkurl{code_cylinder_eq}} \\
Corollary 2.3 &
\texttt{\nolinkurl{PaperECompletion.all_frequency_limits_iff}} \\
Example 2.2, full table &
\texttt{\nolinkurl{PaperECompletion.example22_code_table}} \\
Terminating rational codes &
\texttt{\nolinkurl{PaperECompletion.terminating_code_positive_ternary_denominator}} \\
Lemma 3.1 & \texttt{\nolinkurl{CollatzPadic.iterate_le_of_code_eq}} \\
Theorem 3.2 & \texttt{\nolinkurl{CollatzPadic.periodic_bridge}} \\
Theorem 4.1, box to orbit &
\texttt{\nolinkurl{PaperEModularReturn.modular_return_of_box}} \\
Denominator and growth &
\texttt{\nolinkurl{PaperEModularReturn.runCode_den}},
\texttt{\nolinkurl{runCode_den_gt}} \\
Theorem 4.1, unconditional &
\texttt{\nolinkurl{PaperERecurrence.theorem41}},
\texttt{\nolinkurl{box_recurrence}} \\
Mixed-power cancellation &
\texttt{\nolinkurl{PowerPhaseAsymptotics.tendsto_distinct_noninteger_power_average}} \\
Fourier-to-box recurrence &
\texttt{\nolinkurl{FourierBoxRecurrence.exists_ge_fract_box_of_phase}};
\texttt{\nolinkurl{PowerBoxRecurrence.exists_ge_power_fract_box}} \\
Half-open box frequencies &
\texttt{\nolinkurl{FourierBoxCounting.tendsto_fract_box_count}};
\texttt{\nolinkurl{PowerBoxCounting.tendsto_power_fract_box_count}} \\
Corollary 4.2 &
\texttt{\nolinkurl{PaperECorollaries.return_starts_asymptotic}},
\texttt{\nolinkurl{returnStarts_actual}},
\texttt{\nolinkurl{return_in_multiplicative_interval}} \\
Corollary 4.3 &
\texttt{\nolinkurl{PaperECorollaries.signature_parameter_density}},
\texttt{\nolinkurl{modular_return_of_signature}},
\texttt{\nolinkurl{signature_returns_infinite}} \\
Theorem 4.4, exact predicate and errors &
\texttt{\nolinkurl{OOEEffectiveReturn.return_parameter_iff}},
\texttt{\nolinkurl{count_error}},
\texttt{\nolinkurl{error_power_bound}} \\
Theorem 4.4, positive count and witness &
\texttt{\nolinkurl{OOEEffectiveReturn.count_at_witnessCutoff}},
\texttt{\nolinkurl{exists_bounded_modular_return}} \\
Quantitative analytic inputs &
\texttt{\nolinkurl{OOEEffectiveModes.normalized_mode_bound}};
\texttt{\nolinkurl{FejerBox.finite_box_discrepancy}} \\
Lemma 5.2 & \texttt{\nolinkurl{PreimageGrid.count_odd}},
\texttt{\nolinkurl{count_doubled_odd}},
\texttt{\nolinkurl{count_four}} \\
Lemma 5.3 &
\texttt{\nolinkurl{PreimageDomain.closed_domain_for_target}} \\
Lemma 5.4 & \texttt{\nolinkurl{PreimageGrowth.growth_root}} \\
Theorem 5.1 & \texttt{\nolinkurl{PreimageCertificate12.density_21_25}},
\texttt{\nolinkurl{ancestor_density_21_25}} \\
Real exponent \(21/25\) &
\texttt{\nolinkurl{PaperECompletion.ancestor_density_real}} \\
Theorem 6.1 &
\texttt{\nolinkurl{PreimageBalance.certificate_power_ceiling}},
\texttt{\nolinkurl{rate_lt_of_rows}},
\texttt{\nolinkurl{rate_lt_of_plus_rows}} \\
Logarithmic ceiling &
\texttt{\nolinkurl{PaperECompletion.certificate_log_ceiling}} \\
Proposition 7.1 &
\texttt{\nolinkurl{BackwardMass.backward_mass_counterexample}} \\
Section 7.1, sublinear count and finite mass &
\texttt{\nolinkurl{SublinearCountingMass.sublinear_counting_finite_mass}} \\
Finite complete trees &
\texttt{\nolinkurl{PaperECompletion.full_prefix_tree_masses}} \\
Proposition 7.2 &
\texttt{\nolinkurl{CollatzMoments.complete_family_moment_loss}},
\texttt{\nolinkurl{PaperECompletion.stopping_word_masses}} \\
Error bound (7.4) &
\texttt{\nolinkurl{FibreMassError.predecessor_mass_error}} \\
Theorem 7.3, deficient rows &
\texttt{\nolinkurl{FibreMass.finite_weight_obstruction}} \\
Theorem 7.3, actual deficient targets &
\texttt{\nolinkurl{FibreDeficit.deficient_mass_not_summable}},
\texttt{\nolinkurl{exists_predecessor_mass_lt}} \\
Theorem 7.3, bounded stopping &
\texttt{\nolinkurl{FibreStopping.bounded_policy_deficit}} \\
Proposition 7.4, plus &
\texttt{\nolinkurl{FibreGeneration.ancestor_reciprocals_not_summable}} \\
Proposition 7.4, minus &
\texttt{\nolinkurl{UniformFibreDistortion.ancestor_reciprocals_not_summable}} \\
Proposition 7.5 &
\texttt{\nolinkurl{FibreSignCoupling.oddReturn_minus_two_mul_add_one}},
\texttt{\nolinkurl{oddReturn_plus_two_mul_sub_one}},
\texttt{\nolinkurl{paired_depth_two_deficit}},
\texttt{\nolinkurl{paired_twoStepMass_progression}} \\*
\end{longtable}

\endgroup

AxiomCheckJugglerCollatzPaper.lean prints the dependencies of the 70
selected declarations. The permitted logical dependencies are propext,
Classical.choice, and Quot.sound. No additional logical axiom or
native-evaluation trust extension belongs to this paper's selected
theorem audit. Theorem 4.1's simultaneous-box premise is discharged by
the included analytic proofs. Review of the correspondence between those
statements and the prose is a separate responsibility.

\section{Appendix C. Quantitative OOE
proof}\label{appendix-c.-quantitative-ooe-proof}

This appendix proves Theorem 4.4. Write \(e(x)=\exp(2\pi i x)\),
\(s=1+2Mt\), \(u=\lfloor s^{9/2}\rfloor\), and
\(v=\lfloor s^{9/4}\rfloor\). The selected 70-declaration audit includes
its complete counting and witness conclusions. The written route below
uses {[}AR24{]}; the formal route derives its own sufficient derivative
tests.

\subsection{C.1 Formal proof route}\label{c.1-formal-proof-route}

For a sum of N consecutive samples, the finite third-derivative test,
with derivative magnitude between lambda and 4 lambda, gives \[
12N\max\{\lambda^{1/6},(N/2)^{-1/2},(N\sqrt\lambda)^{-1/2}\}.
\] The fifth-derivative test, with ratio 6, gives \[
7N\max\{\lambda^{1/30},(N/6)^{-1/8},(N\sqrt\lambda)^{-1/8}\}.
\] These are proved by finite differencing from a second-derivative
estimate, for either constant derivative sign. On real dyadic intervals
the support extends one lattice step beyond the last sample; the actual
OOE derivative chains satisfy the bounds throughout that support. The
module OOEEffectiveModes verifies both frequency axes and obtains the
dyadic constant 32 and all-length constant 128 used below.

FejerBox proves the pointwise finite half-open box estimate, including
saturated arcs and boundary hits. OOEEffectiveReturn identifies that box
with the exact floor residues, chooses the rounded cutoff, removes at
most eight initial parameters, absorbs the logarithm, and extracts the
strict witness bound. Its separate audit covers all 23 theorems. Thus no
analytic cancellation premise remains in Theorem 4.4's Lean statement.
This does not assert formalization of the precise external estimate
(C.5).

\subsection{C.2 Explicit analytic
input}\label{c.2-explicit-analytic-input}

Write \(e(x)=\exp(2\cdot \pi\cdot i\cdot x).\) We use J. Arias de Reyna
{[}AR24{]}, Theorem 11 and Table 1, only at derivative orders \(k=3\)
and \(k=5.\) In notation adapted to this proof, if \(floor(Y)>k,\) f has
continuous derivatives through order k on (X,X+Y{]}, and its k-th
derivative lies between positive lambda and Lambda, then, with
\(D=2^k,\) \[
\frac1Y\left|\sum_{X<m\le X+Y}e(f(m))\right|
\le 11\max\left\{
 \left(\frac{\Lambda}{\lambda Y}\right)^{2/D},
 \left(\frac{\Lambda^2}{\lambda}\right)^{1/(D-2)},
 (\lambda Y^k)^{-2/D}\right\}. \tag{C.5}
\] The source gives constants smaller than 11 for both orders. A
derivative of constant negative sign is handled by replacing f by -f,
which conjugates the exponential sum. No monotonicity hypothesis is
needed in this higher-derivative estimate. This cited result is a
written proof input, not an imported Lean theorem or a new claim of this
project. Our phases are smooth for \(x>0,\) and every retained interval
has \(Y=N\ge 6,\) so \(floor(Y)>5\ge k.\) The normalization is by the
real interval length Y, not by the number of integer summands.

\subsection{C.3 Uniform estimates for every truncated Fourier
mode}\label{c.3-uniform-estimates-for-every-truncated-fourier-mode}

Let \(H\ge 1\) be an integer and let integers \(h_1,h_2\) satisfy
\(0<\max(|h_1|,|h_2|)\le H.\) Consider \[
f_h(x)=\frac{h_1}{2}(1+2Mx)^{9/2}
       +\frac{h_2}{2M}(1+2Mx)^{9/4}. \tag{C.6}
\] On a real interval \(N<x\le 2N\) with \(N\ge \max(H^2,6),\) \[
2MN\le1+2Mx\le5MN. \tag{C.7}
\]

\subsection{Nonzero high-power
coefficient}\label{nonzero-high-power-coefficient}

If \(h_1\) is nonzero, the fifth derivative is \[
f_h^{(5)}(x)=\frac{945}{2}h_1M^5s^{-1/2}
             +\frac{945}{64}h_2M^4s^{-11/4},\qquad s=1+2Mx.
\] The absolute ratio of the second term to the first is at most
\(H\cdot s^{-9/4}/(32M)\le 1/32,\) using \(|h_1|\ge 1,\) \(N\ge H,\) and
\(s\ge 2MN.\) The derivative therefore has the sign of \(h_1\)
throughout the interval. The weaker factors 1/2 and 3/2 around its
leading term give valid bounds \[
\lambda=100|h_1|M^{9/2}N^{-1/2},\qquad
\Lambda=600|h_1|M^{9/2}N^{-1/2}. \tag{C.8}
\] Indeed \(945/(4\cdot sqrt(5))>100\) and
\(2835/(4\cdot sqrt(2))<600.\) Substitution into (C.5), with \(k=5,\)
gives \[
\frac1N\left|\sum_{N<m\le2N}e(f_h(m))\right|
\le32H^{1/30}M^{3/20}N^{-1/60}. \tag{C.9}
\] For clarity, the three terms before enlargement are bounded by
\(11\cdot (6/N)^{1/16},\)
\(11\cdot (3600\cdot H\cdot M^{9/2}\cdot N^{-1/2})^{1/30},\) and
\(11\cdot (100\cdot M^{9/2}\cdot N^{9/2})^{-1/16}.\) The scalar
inequalities \(6<2^16\) and \(3600<2^30\) suffice for the constant 32.

\subsection{Vanishing high-power
coefficient}\label{vanishing-high-power-coefficient}

If \(h_1=0,\) then \(h_2\) is nonzero and \[
f_h^{(3)}(x)=\frac{45}{16}h_2M^2s^{-3/4}.
\] Use \[
\lambda=\tfrac12|h_2|M^{5/4}N^{-3/4},\qquad
\Lambda=2|h_2|M^{5/4}N^{-3/4}.
\] The endpoint comparisons follow from \(45/(16\cdot 5^{3/4})>1/2\) and
\(45/(16\cdot 2^{3/4})<2.\) Estimate (C.5), now with \(k=3,\) bounds the
normalized sum by \(22\cdot H^{1/6}\cdot M^{5/24}\cdot N^{-1/8}.\) The
three terms before enlargement are at most \(11\cdot (4/N)^{1/4},\)
\(11\cdot (8\cdot H\cdot M^{5/4}\cdot N^{-3/4})^{1/6},\) and
\(11\cdot ((1/2)\cdot M^{5/4}\cdot N^{9/4})^{-1/4}.\) Each is bounded by
the displayed common expression using \(H,M,N\ge 1.\) Since
\(H^2\le N,\) \[
H^{1/6}N^{-1/8}
\le H^{1/30}N^{-7/120}
\le H^{1/30}N^{-1/60}.
\] Together with (C.9), this proves the single bound \[
\frac1N\left|\sum_{N<m\le2N}e(f_h(m))\right|
\le32H^{1/30}M^{1/4}N^{-1/60}. \tag{C.10}
\] All signs and both coordinate-axis cases are covered.

\subsection{Passage to an initial
segment}\label{passage-to-an-initial-segment}

Suppose \(1\le H\le T^{1/4}.\) Split (0,T{]} into intervals (N,2N{]}
with \(N=T/2,\) T/4, \ldots{} while \(N\ge \max(H^2,6).\) The remaining
initial interval is (0,R{]}, where \(R<2\cdot \max(H^2,6);\) take
\(R=T\) if no interval is retained. The intervals are disjoint as
half-open sets, even at noninteger endpoints. The remainder has floor(R)
integer summands. The change from indices 1,\ldots,T to 0,\ldots,T-1 is
exactly \(e(f_h(0))-e(f_h(T)),\) of norm at most 2. Their total
contribution is at most \(floor(R)+2\le 2H^2+14\le 16H^2,\) including
when no interval is retained. Moreover, \[
\sum_{j\ge0}(T/2^{j+1})^{59/60}
=\frac{T^{59/60}}{2^{59/60}-1}<2T^{59/60}.
\] Since \(H^2\le T^{1/2},\) (C.10) and the initial remainder imply \[
\frac1T\left|\sum_{0\le t<T}e(f_h(t))\right|
\le128M^{1/4}H^{1/30}T^{-1/60}. \tag{C.11}
\] The intermediate constant is at most \(64+16=80;\) 128 is a
convenient enlargement. Integer endpoints require no integrality of the
dyadic N.

\subsection{C.4 An explicit half-open box
estimate}\label{c.4-an-explicit-half-open-box-estimate}

For arbitrary points \(z_0,...,z_{T-1}\) on the two-dimensional unit
torus, let \(E_H\) bound the normalized sums of every nonzero Fourier
mode with \(\max(|h_1|,|h_2|)\le H.\) For every half-open product of
circular intervals B, \[
\left|\frac{\#\{t<T:z_t\in B\}}T-|B|\right|
\le\frac5{\sqrt{H+1}}+(3+2\log H)^2E_H. \tag{C.12}
\] Here is a proof with constants and endpoints retained. The Fejer
kernel \[
F_H(x)=\frac1{H+1}\left|\sum_{j=0}^H e(jx)\right|^2
\] is nonnegative, has integral one, and has Fourier coefficients
1-\textbar h\textbar/(H+1) for \(|h|\le H\) and zero otherwise. For
\(0<\delta\le 1/2,\) \(\sin(\pi\cdot x)\ge 2\cdot x\) on
\(0\le x\le 1/2\) gives \[
\int_{\|x\|\ge\delta}F_H(x)\,dx
\le\frac1{2(H+1)\delta}.
\] The product kernel \(K=F_H\) tensor \(F_H\) therefore has mass at
most \(q=1/((H+1)\cdot \delta)\) outside the coordinate
delta-neighbourhood of zero. Expand and contract each interval of B by
delta on the circle, obtaining \(B^+\) and \(B^-.\) A contracted
interval can be empty and an expanded one full; both cases are retained.
More precisely, write a proper circular interval as the image of
{[}a,a+ell), with \(0<ell<1.\) Its expansion is the image of
{[}a-delta,a+ell+delta) if \(ell+2\cdot \delta<1,\) and the whole circle
otherwise. Its contraction is the image of {[}a+delta,a+ell-delta) if
\(ell>2\cdot \delta,\) and empty otherwise. Empty and full original
intervals are left unchanged. For shifts y with each circular distance
\(||y_i||<\delta,\) these definitions give z in B implies z-y in
\(B^+,\) and z-y in \(B^-\) implies z in B. The exceptional shifts have
K-mass at most q. Since \(0\le 1_B\cdot K\le 1,\) the following
inequalities hold at every point, including the original interval
endpoints: \[
\mathbf1_{B^-}*K-q\le\mathbf1_B\le\mathbf1_{B^+}*K+q.
\] The changes in volume are at most \(4\cdot \delta.\) A circular
interval's Fourier coefficient has absolute value at most
1/max(1,\textbar h\textbar), also for the empty or full interval. Thus
the sum of absolute nonconstant coefficients of either convolution is at
most \[
\left(1+2\sum_{h=1}^H\frac1h\right)^2-1
\le(3+2\log H)^2.
\] Averaging the pointwise inequalities proves an error at most
\(4\cdot \delta+q+(3+2\cdot \log(H))^2\cdot E_H.\) Take
\(\delta=(H+1)^{-1/2}\) for \(H\ge 3.\) For \(H=1\) or 2, the trivial
discrepancy bound one already implies (C.12).

\subsection{C.5 Assembly and first-witness
extraction}\label{c.5-assembly-and-first-witness-extraction}

Apply (C.12) to \[
z_t=\left(\frac{s^{9/2}}2,\frac{s^{9/4}}{2M}\right)\pmod1,
\qquad B=[0,1/2)\mathbin{\times}[1/(2M),1/M).
\] Its area is 1/(4M). For any \(x\ge 0\) and integers \(m\ge 1\) and
\(0\le r<m,\) writing \(floor(x)=m\cdot q+r'\) with \(0\le r'<m\) gives
\(fract(x/m)=(r'+fract(x))/m.\) Consequently \(floor(x)=r\) modulo m is
equivalent to \(r/m\le fract(x/m)<(r+1)/m,\) including equality at the
left endpoint. Box membership is therefore exactly u even and \(v=1\)
modulo 2M. For \(M=1\) the second interval is {[}1/2,1): its right
endpoint is excluded; an integer value of x/(2M) has fractional part
zero and is rejected. The first interval likewise excludes its endpoint
1/2. No absence of boundary hits is assumed. Choose
\(H=floor(T^{1/32}),\) so \(H\ge 1,\) \(H\le T^{1/4},\) and \[
(H+1)^{-1/2}\le T^{-1/64},\qquad
H^{1/30}T^{-1/60}\le T^{-1/64},\qquad
2\log H\le\frac{\log T}{16}.
\] Equations (C.11) and (C.12) give the first bound in (4.7) before the
final +8. Exactly min(T,1+floor(7/M)) parameters have \(1+2Mt<16.\)
Deleting the box hits among these parameters changes the count by at
most eight.

For the second bound in (4.7), put \(x=\log(T)\ge 0.\) The maximum of \[
(3+x/16)^2e^{-x/128}
\] on this half-line is \(256\cdot \exp(-13/8),\) attained at \(x=208,\)
and is less than 64. For example, the first four terms of the
exponential series already give \(\exp(13/8)>4.\) Absorb 8/T into
\(8\cdot T^{-1/128}\) and \(5\cdot T^{-1/64}\) into
\(5\cdot T^{-1/128}.\) Since \(M\ge 1\) and
\(5+128\cdot 64+8=8205<2^14,\) the second inequality follows.

At \(T=2^{2176}M^{160},\) one has \(T^{1/128}=2^17\cdot M^{5/4}.\) Hence
the error is at most T/(8M), while the main term is T/(4M). Thus
\(A_M(T)\ge\) \(T/(8M)>0.\) Select a counted t. Since \(t<T\) is
integral, \(1+2Mt<2MT,\) giving the parameter and start bounds in
Theorem 4.4.

The actual orbit is \[
s^2\longmapsto s^3\longmapsto u\longmapsto v.
\] Both initial sources are odd, the third is even by construction, and
the exact identity \(isqrt(isqrt(s^9))=floor(s^{9/4})\) identifies the
exit. The already proved size threshold \(s\ge 16\) gives \(v>s^2;\) the
first two images are larger still. Finally \(s=1\) modulo 2M gives
\(n=1\) modulo 2M, and the box gives the exit congruence in Theorem 4.4.
This completes the written proof.

\clearpage

\section{References}\label{references}

\begingroup\small

{[}T76{]} R. Terras, \emph{A stopping time problem on the positive
integers}, Acta Arithmetica 30 (1976), 241-252.
\href{https://doi.org/10.4064/aa-30-3-241-252}{doi:10.4064/aa-30-3-241-252}.

{[}BL96{]} D. J. Bernstein and J. C. Lagarias, \emph{The 3x+1 conjugacy
map}, Canadian Journal of Mathematics 48 (1996), 1154-1169.
\href{https://doi.org/10.4153/CJM-1996-060-x}{doi:10.4153/CJM-1996-060-x}.

{[}PP25{]} V. Prasad and M. A. Prasad, \emph{Estimates of the maximum
excursion constant and stopping constant of juggler-like sequences},
preprint, January 2025.
\href{https://doi.org/10.13140/RG.2.2.14110.04168}{doi:10.13140/RG.2.2.14110.04168}.

{[}KL03{]} I. Krasikov and J. C. Lagarias, \emph{Bounds for the 3x+1
problem using difference inequalities}, Acta Arithmetica 109 (2003),
237-258. \href{https://arxiv.org/abs/math/0205002}{arXiv:math/0205002}.

{[}S26{]} M. Sharpe, \emph{Collatz}, Lean source repository, 2026;
Grid50.lean and KLGrid.lean, MIT license; accessed 22 September 2026.
\href{https://github.com/msharpe248/collatz}{Source repository}. The
reused material is attributed in the archived local PreimageGrid.lean
and PreimageGrowth.lean.

{[}Bo94{]} M. D. Boshernitzan, \emph{Uniform distribution and Hardy
fields}, Journal d'Analyse Mathematique 62 (1994), 225-240.
\href{https://doi.org/10.1007/BF02835955}{doi:10.1007/BF02835955}.

{[}R26{]} M. Reilly, \emph{A criterion for weighted uniform distribution
along functions from a Hardy field}, preprint, 2026.
\href{https://arxiv.org/abs/2606.08040v1}{arXiv:2606.08040v1}. Theorem
1.1 states the ordinary criterion used in Section 4.

{[}A{]} P. Cochin, \emph{Lower Bounds for Cycle Lengths in the Juggler
Map}, preprint, 2026.
\href{https://doi.org/10.5281/zenodo.22676452}{Concept
DOI:10.5281/zenodo.22676452}.

{[}B{]} P. Cochin, \emph{Five-Step Descent Certificates for the Juggler
Map: Parity Statistics of Nested Floor Powers}, preprint, 2026.
\href{https://doi.org/10.5281/zenodo.22864933}{Concept
DOI:10.5281/zenodo.22864933}.

{[}C{]} P. Cochin, \emph{Fate Contagion and Termination Criteria for the
Juggler Map}, preprint, 2026.
\href{https://doi.org/10.5281/zenodo.22678164}{Concept
DOI:10.5281/zenodo.22678164}.

{[}D{]} P. Cochin, \emph{No m-cycles of the 3n-1 map for m at most 61},
local version 1.1.1, 2026. Deposited version 1.0.0 states 58.
\href{https://doi.org/10.5281/zenodo.22876189}{Concept
DOI:10.5281/zenodo.22876189}.

{[}Z08{]} R. E. Zarnowski, \emph{The congruence structure of the 3x+1
map}, The Fibonacci Quarterly 46/47 (2008/2009), Theorem 3 and Remark 1.
\href{https://www.fq.math.ca/Papers1/46_47-2/Zarnowski_1-09.pdf}{PDF}.

{[}Tao22{]} T. Tao, \emph{Almost all orbits of the Collatz map attain
almost bounded values}, Forum of Mathematics, Pi 10 (2022), e12.
\href{https://doi.org/10.1017/fmp.2022.8}{doi:10.1017/fmp.2022.8}.

{[}AR24{]} J. Arias de Reyna, \emph{Explicit van der Corput's d-th
derivative estimate}, preprint, version 1, 2024, Theorem 11 and Table 1.
\href{https://arxiv.org/abs/2407.02094v1}{arXiv:2407.02094v1}.

\endgroup
\end{document}
